% ------------------------------------------------------------
% Page 73
% ------------------------------------------------------------

\noindent
C’est à cette époque, que la chambre de Ginette fut totalement détruite. Sa couverture
chauffante embrasa tout ce qui s’y trouvait ; heureusement elle n’y était pas et les pompiers
parvinrent à circonscrire l’incendie au couloir attenant, mais bijoux, lettres et photos, glaces et
bibelots, tout avait disparu.

\vspace{0.45cm}

Après la mort accidentelle de son maître d’hôtel, Ginette, quelques mois plus tard, tomba
malade et fut emportée par un cancer. Son mari, désormais seul, se résigna à vendre cette
vieille demeure à laquelle il était tellement attaché. Ce fut, pour lui, un crève-cœur. Il tourna
donc définitivement le dos à Meyrueis, Ayres et son château en 1979. Une maison de retraite
l’accueillit aux Sables d’Olonne où il termina sa vie.

\vspace{0.45cm}

C’est désormais la famille de Monjou qui préside aux destinées du Château d’Ayres\ldots

\vspace{1.6cm}

% ------------------------------------------------------------
% Encadré
% ------------------------------------------------------------

\noindent
\fbox{%
\parbox{0.89\textwidth}{%
\small
\textit{%
*Frédéric H., de père anglais, tailleur pour hommes, avenue Matignon, ancien combattant
14-18, blessé en 1917, milite dans les Unions Chrétiennes de Jeunes Gens (UCJG).
Sous le couvert de ses activités commerciales et de ses engagements humanitaires où il
soutient des maisons d’enfants, cas sociaux, handicapés\ldots, il peut assez facilement passer la
ligne de démarcation entre la Zone Occupée à la Zone non Occupée et de là se rendre en
Suisse où il s’engage dans un réseau de résistance, le réseau Gilbert, organisé de Genève par
le Colonel Georges Groussard.
Ce Réseau dispose de 700 agents en relation directe avec les services secrets britanniques.
Parmi eux de nombreux « unionistes » Maurice Costil, Charlie Vernier, Le Gal Basteau,
Charles Guillon, Secrétaire International des YMCA (Young Men Christian Association) qui
sont tous plus ou moins convoyeurs de fonds destinés à l’aide aux réfugiés au Chambon sur
Lignon, pour la Cimade etc.%
}
}%
}

\vspace{0.15cm}

\begin{flushright}
Jacques Poujol : « Protestants dans la France en Guerre » p. 228-230
\end{flushright}

\vfill

